% SLOT_ID: the 18 bins of full scale and where each code reads
% (rfcs/0001-g3/04-addressing.md#slot_id).
\documentclass[border=6pt]{standalone}
\usepackage{g3-fig}
\begin{document}
\begin{tikzpicture}
  \def\W{16}            % cm for full scale
  \pgfmathsetmacro\b{\W/18} % cm per bin
  \tikzset{bin/.style={draw, minimum height=1.0cm, minimum width=\b cm, inner sep=0pt,
                       anchor=south west}}
  \tikzset{lbl/.style={anchor=north, inner sep=1pt, font=\ttfamily\scriptsize}}
  \colorlet{bad}{black!30}
  \node[bin, fill=bad] at (0, 0) {};
  \node[lbl] at (0.5*\b, 0.97) {short};
  \node[bin, fill=bad] at (17*\b, 0) {};
  \node[lbl] at (17.5*\b, 0.97) {open};
  \foreach \n in {0,...,15} {
    \pgfmathtruncatemacro\lo{\n/4}
    \pgfmathsetmacro\shade{\lo==0 ? 20 : (\lo==1 ? 35 : (\lo==2 ? 20 : 35))}
    \node[bin, fill=g3Bus!\shade!white] at ({(\n+1)*\b}, 0) {};
    \node[lbl] at ({(\n+1.5)*\b}, 0.97) {\n};
  }
  % nominal readings (ratio R_n / (12.7 k + R_n))
  \foreach \r in {0.0830,0.1390,0.1957,0.2494,0.3068,0.3602,0.4172,0.4708,
                  0.5296,0.5836,0.6402,0.6932,0.7510,0.8046,0.8611,0.9168}
    \draw[very thick, g3Stop] (\r*\W, 0.05) -- (\r*\W, 0.5);

  % scale
  \foreach \f/\t in {0/0, 0.25/0.25, 0.5/0.5, 0.75/0.75, 1/1}
    \draw (\f*\W, 0) -- (\f*\W, -0.3) node[below, note] {\t};
  \node[note] at (\W/2, -1.0) {reading as a fraction of full scale (ADC reference)};

  % loaf braces
  \foreach \l/\a in {0/0x10--0x13, 1/0x14--0x17, 2/0x18--0x1B, 3/0x1C--0x1F}
    \draw[brace] ({(4*\l+1)*\b}, 1.15) -- ({(4*\l+5)*\b}, 1.15)
      node[midway, above=4pt, note] {loaf \l: \texttt{\a}};

  \node[note, anchor=north west, align=left, text width=\W cm] at (0, -1.4)
    {Code \emph{n} = 4 $\times$ loaf + slot reads in bin \emph{n}+1; address 0x10 + \emph{n}.
     Red marks, in the lower half of each bin, are the nominal readings of the E96 values with a 12.7\,k pull-up.
     The worst case with 1\,\% parts and 0--10\,$\Omega$ of switch resistance stays
     2.08\,\% of full scale inside its bin, before ADC error (estimate).};
\end{tikzpicture}
\end{document}
